% ============================================================================
% Xin Yao -- academic curriculum vitae
% Build:  latexmk -pdf Xin_Yao_CV.tex      (or: pdflatex Xin_Yao_CV.tex twice)
% Then copy Xin_Yao_CV.pdf to ../assets/Xin_Yao_CV.pdf and push.
% ============================================================================
\documentclass[11pt,letterpaper]{article}

\usepackage[margin=0.8in,top=0.7in,bottom=0.7in]{geometry}
\usepackage[utf8]{inputenc}
\usepackage[T1]{fontenc}
\usepackage{newtxtext,newtxmath}
\usepackage{microtype}
\usepackage{enumitem}
\usepackage{titlesec}
\usepackage{xcolor}
\usepackage{tabularx}
\usepackage[hidelinks]{hyperref}
\usepackage{fancyhdr}
\usepackage{lastpage}

% ---------- section style: small caps + rule (matches the previous CV) -----
\titleformat{\section}{\large\scshape\bfseries}{}{0em}{}[\vspace{-0.55em}\rule{\textwidth}{0.6pt}]
\titlespacing*{\section}{0pt}{1.1em}{0.5em}
\titleformat{\subsection}{\normalsize\bfseries}{}{0em}{}
\titlespacing*{\subsection}{0pt}{0.7em}{0.25em}

% ---------- lists ------------------------------------------------------------
\setlist[itemize]{leftmargin=1.2em,itemsep=0.15em,topsep=0.2em,parsep=0pt,label=\textbullet}
\setlist[enumerate]{leftmargin=2.6em,itemsep=0.35em,topsep=0.2em,parsep=0pt}

% ---------- helpers ----------------------------------------------------------
% \entry{left bold}{right-aligned date}
\newcommand{\entry}[2]{\noindent\textbf{#1}\hfill\textbf{#2}\par}
% \sub{plain line under an entry}
\newcommand{\sub}[1]{\noindent #1\par}
% \proj{project title}
\newcommand{\proj}[1]{\vspace{0.45em}\noindent\textbf{#1}\par\vspace{0.1em}}
% \meta{institution / advisor / related publications}
\newcommand{\meta}[1]{\noindent{\small\textit{#1}}\par\vspace{0.1em}}
% \me = my own name in author lists
\newcommand{\me}{\textbf{X. Yao}}

\pagestyle{fancy}
\fancyhf{}
\renewcommand{\headrulewidth}{0pt}
\fancyfoot[C]{\small Xin Yao \textperiodcentered\ Curriculum Vitae \textperiodcentered\ \thepage\ of \pageref{LastPage}}

\setlength{\parindent}{0pt}
\setlength{\parskip}{0.2em}

\begin{document}

% ============================================================================
\begin{center}
  {\LARGE\scshape Xin Yao}\\[0.35em]
  1928 SW 42nd Drive, Gainesville, FL 32607, USA\\
  (+1) 571-236-1806 \textperiodcentered\ \href{mailto:xyao@ufl.edu}{xyao@ufl.edu}
  \textperiodcentered\ \href{https://xyao2021.github.io/homepage/}{xyao2021.github.io/homepage}\\
  \href{https://scholar.google.com/citations?user=gwTaRUkAAAAJ&hl=en}{Google Scholar}
  \textperiodcentered\ \href{https://www.linkedin.com/in/xin-yao-aa9249204/}{LinkedIn}
  \textperiodcentered\ \href{https://github.com/XYao2021}{GitHub}
\end{center}

% ============================================================================
\section{Research Interests}
My research lies at the intersection of machine learning, wireless communication systems, and
information theory. I combine theoretical analysis with algorithm and system design to develop
efficient, secure, and adaptive methods for distributed learning and dynamic network control, and
evaluate them on real-world wireless testbeds that I build from the physical layer up.

\textbf{Keywords:} reinforcement learning for wireless networks; Lyapunov optimization under
real-world constraints; LoRa and low-power multi-hop networks; software-defined radio and
physical-layer implementation; wireless testbeds and remote experimentation platforms; decentralized
and federated learning; error feedback and communication compression; information-theoretic secure
aggregation; mmWave and RIS-assisted systems; wireless sensing and indoor localization.

% ============================================================================
\section{Research Statement (Summary)}
The long-term goal of my research is to make learning-driven wireless networks practical: networks that use
machine learning to allocate resources, share spectrum, and adapt the physical layer, with guarantees of
convergence, privacy, and reliability that continue to hold after deployment on real hardware. A central
motivation is the gap between simulation and reality: many learning algorithms proposed for wireless systems
perform well in simulation but degrade on real hardware, where carrier-frequency offsets, timing errors,
interference, and hardware outages violate their design assumptions. I address this gap from both directions,
by developing algorithms whose analysis accounts for real-world constraints and by building experimental
infrastructure that makes over-the-air validation systematic.

My doctoral research follows three threads: (1) \emph{adaptive error feedback for distributed learning},
where A-DEFEAT compensates for communication-compression error in decentralized learning with a new proof
technique and convergence analysis; (2) \emph{reinforcement learning for real-world wireless networks}, where
objectives are derived from Lyapunov functions that encode real-world constraints, state representations are
built from quantities a deployed system can observe, and Q-learning, DDPG, and MADDPG policies are compared in
simulation before deployment (spectrum sharing on a real-time LoRa testbed, LDPP-MADDPG power allocation, and
FORGE-LoRa multi-hop adaptation); and (3) \emph{wireless testbed development}, including the physical-layer
stack of the EdgeAI Lab testbed, the UnionLabs remote-access platform, and UAV nodes.

Future plans build on these results in three directions: cross-layer learning-based adaptation for
heterogeneous networks with safe online learning under non-stationary conditions; digital-twin pipelines
calibrated from testbed measurements, with quantitative criteria for when a policy trained in the twin can be
deployed over the air; and decentralized learning over unreliable wireless networks with time-varying
topologies, lossy links, and quantized communication, connected to secure aggregation under user dropouts. In
parallel, I plan to develop UnionLabs into shared community infrastructure federated across institutions, with
support sought from NSF CNS/NeTS and CCF programs, a CAREER proposal, and industry partnerships in AI-for-RAN.

% ============================================================================
\section{Education}
\entry{University of Florida, Gainesville, FL, USA}{08/2024 -- present}
\sub{Ph.D. Candidate, Electrical and Computer Engineering}
\sub{Advisor: Prof. Mingyue Ji (EdgeAI Lab)}
\sub{Research: machine learning for real-world wireless systems; learning-driven network control on live
testbeds; decentralized learning}

\vspace{0.4em}
\entry{University of Utah, Salt Lake City, UT, USA}{08/2021 -- 07/2024}
\sub{Ph.D. Studies, Electrical and Computer Engineering (transferred to the University of Florida with my
advisor's research group in Aug.\ 2024)}
\sub{Advisor: Prof. Mingyue Ji}

\vspace{0.4em}
\entry{George Washington University, Washington, DC, USA}{01/2020 -- 06/2021}
\sub{M.S., Electrical Engineering}

\vspace{0.4em}
\entry{Tianjin University of Science and Technology (TUST), Tianjin, China}{09/2015 -- 06/2019}
\sub{B.Eng., Electronic Information Engineering}

% ============================================================================
\section{Research Experience}
\entry{Doctoral Researcher, EdgeAI Lab (Advisor: Prof. Mingyue Ji)}{08/2021 -- present}
\sub{University of Utah (08/2021 -- 07/2024) and University of Florida (08/2024 -- present)}

% ----------------------------------------------------------------------------
\proj{EdgeAI Lab: Real-World Wireless Communication Testbed}
\meta{University of Florida \textperiodcentered\ Platform for the over-the-air experiments in [C1], [C4]}
\begin{itemize}
  \item Implemented the physical-layer software stack of the lab's wireless testbed in C++: modulation and demodulation, filtering, synchronization, channel encoding and decoding, and correction of carrier-frequency offset and phase offset.
  \item Built Python APIs on top of the C++ core so that over-the-air experiments can be scripted, repeated,
        and parameterized from a single configuration; designed automated experiment workflows that adapt to
        heterogeneous physical-layer hardware (USRP software-defined radios and LoRa radios).
  \item Deployed the group's learning algorithms (reinforcement-learning spectrum sharing, power allocation,
        and per-slot LoRa PHY adaptation) on both lab-scale and building-scale deployments of the testbed,
        and used the platform for the real-time LoRa experiments reported in [C4] and the multi-hop LoRa
        experiments in [C1].
  \item Technologies: C++, Python, USRP software-defined radios, LoRa radios, Linux.
\end{itemize}

% ----------------------------------------------------------------------------
\proj{UnionLabs: AWS-Based Remote Access and Sharing of NextG and IoT Testbeds}
\meta{University of Florida \textperiodcentered\ NSF-funded project (CIRC award, 2025--2027) \textperiodcentered\ Poster [P2], demo [P3]}
\begin{itemize}
  \item Developed and deployed the AWS-hosted web user interface for remote access to the wireless testbeds.
  \item Designed the physical-layer (PHY) code for different hardware, including USRP software-defined radios and LoRa devices.
  \item Designed the unified \emph{Union API} that reads and drives the different PHY implementations and
        connects them automatically to users' learning or optimization algorithms, and an abstraction layer
        that lets users run the PHY and Union API with a single command.
  \item Supported collaborators from several universities in running their learning and optimization
        algorithms through UnionLabs, including experiments that use multiple testbeds at the same time.
  \item The UnionLabs and OSWireless demonstration won \emph{Third Prize} in the Demo Competition of the
        Warren B.\ Nelms Annual IoT Conference (Dec.\ 2025).
  \item Technologies: AWS, Java, JavaScript, JSON, Python, USRP, LoRa.
\end{itemize}

% ----------------------------------------------------------------------------
\proj{Constrained Optimization for Real-World Wireless Problems with Reinforcement Learning}
\meta{University of Utah and University of Florida \textperiodcentered\ Publications [C4], [C5]; poster [P4]}
\begin{itemize}
  \item Derived optimization objectives with Lyapunov drift-plus-penalty (LDPP) functions under real-world constraints, turning long-term constrained problems into per-slot decisions that a learning agent can act on.
  \item Designed novel state representations for reinforcement-learning algorithms applied to
        communication-system control, and evaluated Q-learning, DDPG, and MADDPG policies both in simulation
        and in deployment on the EdgeAI Lab testbed.
  \item Proposed LDPP-MADDPG, a distributed power-allocation scheme for mmWave cellular networks that combines a Lyapunov drift-plus-penalty objective with multi-agent DDPG [C5]; an earlier MADDPG power-allocation study used measurement data from the POWDER platform [P4].
  \item Carried out an experimental evaluation of reinforcement learning for spectrum sharing on a real-time
        LoRa testbed, comparing Q-learning, DDPG, and MADDPG on live hardware [C4].
  \item Technologies: Python, reinforcement learning, Lyapunov optimization.
\end{itemize}

% ----------------------------------------------------------------------------
\proj{FORGE-LoRa: Network Optimization with Low-Power (LoRa) Communication}
\meta{University of Florida \textperiodcentered\ Publication [C1] (IEEE MILCOM 2026, accepted)}
\begin{itemize}
  \item Derived the optimization objective for multi-hop LoRa networks with Lyapunov functions under
        real-world constraints.
  \item Designed the system architecture for jointly adapting the LoRa PHY parameters on a per-slot basis.
  \item Building on the joint-adaptation method, designed a path discovery and optimization system that keeps
        multi-hop communication always live against interference and hardware outages.
  \item Evaluated the dynamic-control policies in simulation and in deployment on the EdgeAI Lab testbed across
        multiple reinforcement-learning algorithms.
  \item Technologies: LoRa, Python, reinforcement learning, Lyapunov optimization.
\end{itemize}

% ----------------------------------------------------------------------------
\proj{LoFiLoc: Device-Free LoRa Fingerprint Indoor Localization with Commodity Radios}
\meta{University of Florida \textperiodcentered\ Publication [C2] (IEEE MILCOM 2026 Workshops, accepted)}
\begin{itemize}
  \item Device-free indoor localization based on fingerprints of LoRa channel measurements collected with commodity LoRa radios, so that a target carrying no device can be localized from changes in the wireless channel.
\end{itemize}

% ----------------------------------------------------------------------------
\proj{Transformer-Enhanced DDPG for RIS-Assisted MU-MISO Systems (collaboration)}
\meta{University of Florida \textperiodcentered\ Publication [C3] (IEEE MILCOM 2026 Workshops, accepted)}
\begin{itemize}
  \item Contributed to a transformer-enhanced deep deterministic policy gradient (DDPG) method for the joint
        design of base-station beamforming and reconfigurable-intelligent-surface (RIS) phase shifts in
        multi-user MISO systems, with H.\ Zhang and O.\ Jiang.
\end{itemize}

% ----------------------------------------------------------------------------
\proj{A-DEFEAT: Adaptive Error Feedback for Decentralized Learning}
\meta{University of Utah and University of Florida \textperiodcentered\ Poster [P1]; manuscript [M1]}
\begin{itemize}
  \item Designed a new decentralized learning algorithm with an adaptive, discounted error-feedback
        (error-compensation) mechanism for communication compression under data heterogeneity.
  \item Proposed the key proof technique for the adaptive error-feedback mechanism, established convergence
        guarantees for the proposed algorithm, and compared its convergence rate with existing algorithms.
  \item Evaluated the algorithm against existing decentralized learning methods under a range of settings. \item Technologies: Python, optimization theory.
\end{itemize}

% ----------------------------------------------------------------------------
\proj{Information-Theoretic Secure Aggregation with Uncoded Groupwise Keys}
\meta{University of Utah \textperiodcentered\ Publications [J1], [C6] \textperiodcentered\ With K.\ Wan, H.\ Sun, M.\ Ji, and G.\ Caire}
\begin{itemize}
  \item Studied the information-theoretic secure aggregation problem of federated learning, in which the
        server must learn only the aggregate of the users' models while tolerating user dropouts, under the
        practical assumption that keys are mutually independent and each key is shared, uncoded, by a group
        of $S$ users.
  \item The work characterizes when uncoded groupwise keys are optimal: if $S > K-U$ (with $K$ users and
        at most $K-U$ dropouts) a new scheme achieves the same optimal communication cost as the best coded-key
        scheme, whereas if $S \le K-U$ uncoded groupwise key sharing is strictly sub-optimal.
  \item Implemented the proposed secure-aggregation scheme on Amazon EC2 and compared it experimentally with
        existing secure-aggregation schemes that use offline key sharing.
\end{itemize}

% ----------------------------------------------------------------------------
\proj{UAV-Based Wireless Communication Platform}
\meta{University of Florida}
\begin{itemize}
  \item Assembled customized drones and integrated wireless communication payloads for airborne networking
        experiments; initialized and tuned the flight controllers; designed and mounted the customized
        communication payloads.
  \item Designed the communication protocol between ground devices and the airborne node.
  \item Technologies: UAV flight controllers, embedded systems, wireless payload integration.
\end{itemize}

% ----------------------------------------------------------------------------
\proj{Master's Projects}
\meta{George Washington University}
\begin{itemize}
  \item \textbf{Age of Information with Energy Harvesting} (advisor: Prof.\ Omur Ozel). Analyzed the age of
        information (AoI) in M/G/1/1 and idealized infinite-buffer queueing models, including variants with
        packet deadlines, and investigated AoI-minimization strategies for the bufferless M/G/1/1 model with
        packet deadlines.
  \item \textbf{Secure File Transfer: Encryption and Decryption for FTP} (advisor: Prof.\ Tian Lan). Designed a
        custom encryption/decryption scheme for FTP on Linux with client-side encryption and server-side
        decryption of commands and data, and verified end-to-end encrypted FTP transmission.
\end{itemize}

% ----------------------------------------------------------------------------
\proj{Undergraduate Projects}
\meta{Tianjin University of Science and Technology}
\begin{itemize}
  \item \textbf{Voice-controlled smart-home system} (advisor: Jing He). Designed STM32F103-based hardware with
        an LD3320 voice-recognition module; wrote firmware linking the microcontroller to speaker, PC, and
        LCD peripherals; tested spoken commands with system feedback.
  \item \textbf{License-plate recognition system} (advisor: Prof.\ Yang An). Built a MATLAB pipeline for image
        capture from surveillance cameras, filtering (Hamming window), and character segmentation and matching.
  \item \textbf{Android music player} (advisor: Prof.\ Ya-Jun Li). Developed a multi-format music player in Java
        using Eclipse and the Android SDK.
  \item \textbf{Digital counter} (advisor: Prof.\ Yu-liang Liu). Designed the counter circuit, selected
        components, and assembled, tested, and debugged the board.
\end{itemize}

% ============================================================================
\section{Publications}
Google Scholar: \href{https://scholar.google.com/citations?user=gwTaRUkAAAAJ&hl=en}{scholar.google.com/citations?user=gwTaRUkAAAAJ}.
My name is shown in bold.

\subsection{Conference Papers}
\begin{enumerate}[label={[C\arabic*]}]
  \item \me, A.\ Bhuyan, and M.\ Ji, ``FORGE-LoRa: Joint per-Slot PHY Adaptation for Multi-Hop LoRa
        Networks,'' in \emph{Proc.\ IEEE Military Communications Conference (MILCOM)}, National Capital
        Region, USA, Oct.\ 2026. (Accepted, to appear.)\\
        {\small Joint per-slot adaptation of LoRa PHY parameters combined with path discovery and optimization
        so that multi-hop LoRa communication stays live under interference and hardware outages; evaluated in
        simulation and on the EdgeAI Lab testbed.}
  \item \me, A.\ Bhuyan, and M.\ Ji, ``LoFiLoc: Device-Free LoRa Fingerprint Indoor Localization with
        Commodity Radios,'' in \emph{Proc.\ IEEE Military Communications Conference (MILCOM) Workshops},
        National Capital Region, USA, Oct.\ 2026. (Accepted, to appear.)
  \item H.\ Zhang, O.\ Jiang, \me, and M.\ Ji, ``Transformer-Enhanced DDPG for Beamforming and Phase Shift
        in RIS-Assisted MU-MISO Systems,'' in \emph{Proc.\ IEEE Military Communications Conference (MILCOM)
        Workshops}, National Capital Region, USA, Oct.\ 2026. (Accepted, to appear.)
  \item \me, X.\ Wang, H.\ Zhang, and M.\ Ji, ``Experimental Evaluation of Reinforcement Learning for Spectrum
        Sharing on a Real-Time LoRa Testbed,'' in \emph{Proc.\ IEEE International Symposium on Spectrum
        Innovation (DySPAN)}, Washington, DC, USA, May 2026, pp.\ 1--8.
        DOI: \href{https://doi.org/10.1109/DySPAN69846.2026.11571139}{10.1109/DySPAN69846.2026.11571139}.\\
        {\small Q-learning, DDPG, and MADDPG spectrum-sharing policies evaluated on live LoRa hardware in real
        time rather than in simulation only.}
  \item \me, A.\ Bhuyan, X.\ Zhang, and M.\ Ji, ``A Novel LDPP-MADDPG Approach for Distributed Power Allocation
        in mmWave Cellular Networks,'' in \emph{Proc.\ IEEE Military Communications Conference (MILCOM), 6G and
        Spectrum Security for Critical Communication (6GSECC) Workshop}, Los Angeles, CA, USA, Oct.\ 2025.\\
        {\small Distributed power allocation for mmWave cellular networks that combines a Lyapunov
        drift-plus-penalty (LDPP) objective with multi-agent deep deterministic policy gradient (MADDPG).}
  \item K.\ Wan, \me, H.\ Sun, M.\ Ji, and G.\ Caire, ``GroupSecAgg: Information Theoretic Secure Aggregation
        with Uncoded Groupwise Keys,'' in \emph{Proc.\ IEEE International Conference on Communications (ICC)},
        Rome, Italy, May--Jun.\ 2023, pp.\ 3890--3895.
        DOI: \href{https://doi.org/10.1109/ICC45041.2023.10279171}{10.1109/ICC45041.2023.10279171}.\\
        {\small Conference version of [J1].}
\end{enumerate}

\subsection{Journal Articles}
\begin{enumerate}[label={[J\arabic*]}]
  \item K.\ Wan, \me, H.\ Sun, M.\ Ji, and G.\ Caire, ``On the Information Theoretic Secure Aggregation with
        Uncoded Groupwise Keys,'' \emph{IEEE Transactions on Information Theory}, vol.\ 70, no.\ 9,
        pp.\ 6596--6619, Sept.\ 2024.
        DOI: \href{https://doi.org/10.1109/TIT.2024.3422087}{10.1109/TIT.2024.3422087}.
        Preprint: \href{https://arxiv.org/abs/2204.11364}{arXiv:2204.11364}.\\
        {\small Characterizes the optimal communication cost of information-theoretic secure aggregation when
        keys are mutually independent and each key is shared by a group of $S$ users: uncoded groupwise keys
        are optimal when $S > K-U$ and strictly sub-optimal when $S \le K-U$. Includes an Amazon EC2
        implementation compared against existing schemes with offline key sharing.}
\end{enumerate}

\subsection{Manuscripts in Preparation or Under Review}
\begin{enumerate}[label={[M\arabic*]}]
  \item \me, S.\ Wang, J.\ Perazzone, K.\ S.\ Chan, and M.\ Ji, ``Compressed Decentralized Learning with
        Error-Feedback under Data Heterogeneity,'' 2026.
\end{enumerate}

% ============================================================================
\section{Presentations, Posters and Demonstrations}
\subsection{Conference Presentations}
\begin{enumerate}[label={[T\arabic*]}]
  \item ``FORGE-LoRa: Joint per-Slot PHY Adaptation for Multi-Hop LoRa Networks,'' IEEE Military
        Communications Conference (MILCOM), National Capital Region, USA, Oct.\ 2026. (Upcoming.)
  \item ``Experimental Evaluation of Reinforcement Learning for Spectrum Sharing on a Real-Time LoRa Testbed,''
        IEEE International Symposium on Spectrum Innovation (DySPAN), Washington, DC, USA, May 2026.
  \item ``A Novel LDPP-MADDPG Approach for Distributed Power Allocation in mmWave Cellular Networks,'' IEEE
        MILCOM 6G and Spectrum Security for Critical Communication (6GSECC) Workshop, Los Angeles, CA, USA,
        Oct.\ 2025.
  \item ``GroupSecAgg: Information Theoretic Secure Aggregation with Uncoded Groupwise Keys,'' IEEE
        International Conference on Communications (ICC), Rome, Italy, Jun.\ 2023.
  \item ``High-Band Frequency Spectrum for 5G'' (team presentation), George Washington University,
        Washington, DC, USA, 2020.
\end{enumerate}

\subsection{Posters and Demonstrations}
\begin{enumerate}[label={[P\arabic*]}]
  \item ``A-DEFEAT: An Adaptively Discounted Error-Feedback Algorithm for Decentralized Learning'' (poster),
        North American School of Information Theory (NASIT), Brigham Young University, Provo, UT, USA,
        Jun.\ 2026.
  \item ``UnionLabs: Shared Access to Heterogeneous Wireless IoT and Wireless OS Testbeds'' (poster), Warren B.\
        Nelms Annual IoT Conference, University of Florida, Gainesville, FL, USA, Dec.\ 2025.
  \item ``OSWireless: Automatic System to Maximize Network Throughputs'' (live demonstration), Warren B.\ Nelms
        Annual IoT Conference, University of Florida, Gainesville, FL, USA, Dec.\ 2025. \emph{Third Prize,
        Demo Competition.}
  \item ``MADDPG-Based Power Allocation Scheme Using Measurement Data from POWDER'' (poster), Warren B.\ Nelms
        Annual IoT Conference, University of Florida, Gainesville, FL, USA, Dec.\ 2024.
\end{enumerate}

% ============================================================================
\section{Teaching Experience}
\begin{itemize}
  \item \textbf{Lab Instructor, EEL 4514 Communication Systems and Components, University of Florida.} Designed
        and verified the lab sessions on signal processing, held and graded the lab sessions, and held office
        hours for course-related questions before exams.
  \item \textbf{Teaching Assistant, ECE 5510 Random Processes, University of Utah.} Assisted in designing the
        lecture slides, graded homework and exams, and held weekly office hours to answer course-related
        questions.
  \item \textbf{Project Instructor, UnionLabs Web Design (undergraduate project), University of Florida.}
        Directed and guided students through the web-design part of the UnionLabs project, from background
        knowledge to design skills including JavaScript, ReactJS, and the communication logic between the web
        interface and the testbeds.
  \item \textbf{Project Instructor, Automatic System Design (undergraduate project), University of Florida.}
        Guided students to understand the structure and code of the telco system, an AI-in-the-loop evaluation
        system, and how to use the code for customized purposes, and introduced hardware-in-the-loop design
        based on the telco system.
  \item \textbf{Testbed Training and Onboarding, UnionLabs.} Onboard collaborators from multiple universities,
        ranging from machine-learning researchers new to radios to communications engineers new to
        reinforcement learning, onto the shared wireless testbeds: prepare documentation for the customized
        system, design worked examples, and structure first sessions so that every user runs a real
        over-the-air experiment quickly.
\end{itemize}

% ============================================================================
\section{Teaching Statement (Summary)}
\textbf{Philosophy.} Students learn engineering most effectively when theory is connected to the physical
world as early and as often as possible. In wireless communications and machine learning this connection is
concrete: a matched filter derived on paper can be observed as a constellation diagram from the student's own
transmitter, a coding theorem can be verified as a measured shift in a bit-error-rate curve, and a learning
algorithm can be run and measured on real machines. My goal is to help students build a solid conceptual
foundation, the practical skill to implement and debug real systems, and the confidence that comes from seeing
their own analysis confirmed by measurement.

\textbf{Methods.} Theoretical foundations in class, with emphasis on the mathematical expressions; live
software-defined-radio demonstrations so that effects such as intersymbol interference and coding gain are
observed as they are derived; project-based homework tested over the air or against recorded real signals
through remote testbed access of the kind UnionLabs provides; and projects scaffolded in explicit stages from
simulation to hardware, so that students with different preparation all reach the final over-the-air
milestone. Learning is assessed through staged milestones that reward implementation, measurement, and
debugging alongside analysis, and the course is adjusted using student feedback throughout the semester.

\textbf{Mentoring.} Early contact with the real system followed by growing ownership: a new student runs a
working over-the-air experiment in the first weeks, then takes responsibility for one testbed component or
algorithm, growing from reproducing a known result to extending it to leading a first-authored paper, with
regular individual meetings, projects matched to each student's strengths, and frequent presentations.

\textbf{Courses I am prepared to teach.} Signals and systems, digital communications, and digital signal
processing (undergraduate); wireless system design and networking courses that connect the physical and MAC
layers; machine learning for wireless networks (graduate); and a new project-based course, \emph{Hands-On
Wireless Systems with Software-Defined Radios}, in which students implement a complete communication link and
then apply a learning algorithm to control it.

% ============================================================================
\section{Technical Skills}
\begin{itemize}
  \item \textbf{Programming:} Python, C/C++, Java, JavaScript, MATLAB, Arduino, JSON, XHTML/DHTML.
  \item \textbf{Machine learning:} reinforcement learning (Q-learning, DDPG, MADDPG, transformer-enhanced
        policies), federated learning, decentralized learning, communication compression and error feedback,
        secure aggregation, distributed and parallel computing.
  \item \textbf{Wireless communications and signal processing:} physical-layer implementation (modulation and
        demodulation, synchronization, filtering, channel encoding and decoding, carrier-frequency and
        phase-offset correction), software-defined radio (USRP), LoRa and low-power multi-hop networking,
        mmWave power allocation, RIS-assisted systems, real-world testbed development, network protocol
        analysis, wireless fingerprinting for localization.
  \item \textbf{Optimization and theory:} Lyapunov drift-plus-penalty optimization, convergence analysis of
        decentralized learning, information-theoretic secure aggregation, queueing models and age of
        information.
  \item \textbf{Embedded, cloud, and systems:} STM32 microcontrollers, Android SDK (Eclipse/JDK), AWS (EC2 and
        web deployment), Linux, UAV flight-controller setup and payload integration.
  \item \textbf{Design and productivity tools:} AutoCAD, Siemens NX, Sharp3D, Adobe (Photoshop, Premiere Pro,
        After Effects, Acrobat), MS Office, \LaTeX.
\end{itemize}

% ============================================================================
\section{Honors and Awards}
\begin{itemize}
  \item Third Prize, Demo Competition, Warren B.\ Nelms Annual IoT Conference, University of Florida
        (UnionLabs / OSWireless demonstration) \hfill 2025
  \item Third-Prize Scholarship, Tianjin University of Science and Technology \hfill 2017
  \item Advanced Individual of Social Work, Tianjin University of Science and Technology \hfill 2016
  \item Excellent League Member, Tianjin University of Science and Technology \hfill 2016
\end{itemize}

\vspace{0.6em}
{\small\textit{Last updated: September 2026.}}

\end{document}
